\section{Installation, startup and station identity}\label{sec:scope}
The monitor combines rocket telemetry, Away antenna pointing, local USB video,
recording/replay and a nominal OpenRocket reference. It runs offline for local
devices, saved flights, demos, coordinates and prepared simulations. Internet
access is needed for online weather and initial build downloads. The monitor
does not require internet for its existing direct station LAN connection.

\subsection{Install the released package}
Download the platform asset and optional SHA-256 sidecar from the
\href{https://github.com/MIT-Rocket-Team/RT_Telemetry_UI/releases/tag/v0.1.1}{0.1.1 release}.
On macOS 14 or newer, open the universal DMG and copy the complete application
to Applications. It selects a native Apple Silicon or Intel payload.
The release is ad-hoc signed and not notarized; if macOS blocks first launch,
use the operating system's explicit approval for the trusted downloaded app.
On Windows, extract the entire x64 ZIP and run the executable within its folder.
Keep \code{_internal} and all adjacent files together. Windows ARM runs this
x64 distribution through emulation. Linux is not a validated release target.

Python, Qt, FFmpeg, Java and OpenRocket are included. Recipients do not install
development tools. USB drivers and camera permissions depend on the computer;
Dire Wolf is a separate optional dependency for APRS. Lucida Grande is loaded
privately from the package without installing a system font.

\subsection{The four-page launch wizard}
Choose \ui{Launch station} or \ui{Away 1--4}, then \ui{Telemetry} or
\ui{Video}, \ui{Balius} or \ui{Iris}, and \ui{URRG} or a custom site.
Station identity and role remain fixed for this session. Finish opens LIVE
with serial devices disconnected and cameras stopped. Cancel exits before a
device controller is constructed.

\ui{Remember these choices next time} is off by default. Without it, startup
defaults are Launch / Telemetry / Balius / Custom, subject to CLI overrides.
With it, the wizard is prefilled from the saved choices. Uncheck it and finish
to stop remembering. Cancel preserves the previous remembering preference.
\ui{File $\to$ Forget setup on next launch} clears only startup choices, has
no keyboard shortcut, and does not alter the running mission or global settings.

New missions are named \ui{Balius Launch} or \ui{Iris Launch}. URRG fills the
pad and applicable station coordinates. Elevations remain separate entries.
Opening a saved mission or flight restores that file's name and locations;
startup defaults do not replace them.

\subsection{Command-line selection}
From the standalone source repository:
\begin{lstlisting}[language=bash]
.venv/bin/python -m rocket_gnc_monitor \
  --site away3 --role telemetry --vehicle iris \
  --launch-site urrg --skip-setup
\end{lstlisting}
The packaged executable accepts the same arguments. On Mac it resides inside
\nolinkurl{RocketGNCMonitor-v0a.app/Contents/MacOS/RocketGNCMonitor-v0a}; keep the
complete bundle in place. On Windows invoke the extracted executable.

\begin{longtable}{@{}L{.30\linewidth}L{.63\linewidth}@{}}
\toprule Option & Behavior \\\midrule\endhead
\code{--site} & \code{launch}, \code{away1}, \code{away2}, \code{away3}, \code{away4}. The old \code{base} spelling is a compatibility alias. \\
\code{--role} & \code{telemetry} or \code{video}. \\
\code{--vehicle} & \code{balius} or \code{iris}. \\
\code{--launch-site} & \code{urrg} or \code{custom}. \\
\code{--skip-setup} & Start directly with resolved choices; otherwise review them in the wizard. \\
\code{--remember-setup} & Check the opt-in; with skipped setup, save the resolved choices. \\
\code{--no-remember-setup} & Clear the opt-in; with skipped setup, remove remembered choices. Without either remember flag, skipped setup leaves saved choices unchanged. \\
\code{--station} & Legacy shortcut: \code{launch}, \code{away}, \code{video} preselect Launch Telemetry, Away 1 Telemetry, Away 1 Video. Does not bypass setup. \\
\code{--demo} & Explicitly request recorded DEMO instead of normal LIVE startup. \\
\code{--demo-station} & \code{GS1}, \code{GS2}, \code{GS3}; default GS2. \\
\code{--data-dir} & Separate per-user application data/preferences directory. \\
\bottomrule
\end{longtable}

\section{Workspaces and information placement}\label{sec:workspaces}
Station role and source mode are independent. The same telemetry workspace can
show LIVE, DEMO or REPLAY. None of those source-mode names is the rocket's
onboard state. Video does not occupy either telemetry workspace.

\begin{figure}[H]\centering
\begin{tikzpicture}
\node[box,text width=71mm,minimum height=35mm,anchor=north] (flight) at (0,0) {\textbf{Launch / Telemetry: Flight}\\[4pt]Main telemetry and link state\\Compact rocket + trajectory / 3D\\GNC plots and actuator matrix\\Away communication-route selector};
\node[box,text width=71mm,minimum height=35mm,anchor=north] (systems) at (8.2,0) {\textbf{Launch / Telemetry: Systems}\\[4pt]Full telemetry, GPS, servos, cells\\Commands, power/BMS, recovery\\Mission, wind and OpenRocket\\Recording, diagnostics and network};
\node[accentbox,text width=129mm] (owner) at (4.1,-4.6) {One controller, mission and recording session shared by both Launch displays};
\draw[flow] (owner.north west)--(flight.south);
\draw[flow] (owner.north east)--(systems.south);
\node[box,text width=71mm,minimum height=23mm,anchor=north] at (0,-6) {\textbf{Away / Telemetry: one window}\\Main rocket telemetry and pose\\Local antenna graphic and controls};
\node[box,text width=71mm,minimum height=23mm,anchor=north] at (8.2,-6) {\textbf{Video: separate operator role}\\Away: two assigned local inputs\\Launch: large channels + thumbnails};
\end{tikzpicture}
\caption{Information placement in v0a. These are functional layouts, not pixel-exact screenshots. Configuration and file dialogs remain deliberate actions.}
\label{fig:workspaces}
\end{figure}


\subsection{Launch Telemetry: Flight and Systems}
\ui{Flight} shows primary telemetry, state/freshness, the ground-board and
rocket-link indicators, flight plots/reference, the full 3D viewer, compact
rocket orientation display, GNC plots and actuator channels. Its communication
panel selects a desired route through Away 1--4. Launch has no local physical
or virtual antenna pointer.

\ui{Systems} shows the complete legacy telemetry and GPS tables, servos, cells,
rocket commands, power/BMS and recovery controls, mission/wind/OpenRocket,
sessions/replay, diagnostics and APRS/network. A large state badge occupies the
upper-right corner. Both windows share one controller, mission and recording
session. There are no duplicated serial workers.

The intended layout is two 1920 by 1080 displays, approximately 1920 by 1020
usable window space. \ui{View $\to$ Arrange station displays} positions the
windows on available screens. Both remain accessible on one screen for
preparation. Closing Flight exits the workspace; the Systems close button
keeps that required display open. Long diagnostics, JSON and wind profiles
scroll within visible panels. Smaller windows/display scaling can crowd
controls; this release does not resolve every compact-layout issue.

\subsection{Away Telemetry and Video roles}
Away Telemetry places main rocket telemetry, the large compact rocket and its
altitude plot beside the local articulated pointer and controls. Advanced
mission, settings and file operations remain available through the menus and
dialogs. Its downlink and pointer boards connect independently.

Away Video presents the two USB inputs assigned to that station. Launch Video
presents two Balius or three Iris large channels with eight Away thumbnails
below. Video roles hide serial controls and rocket instruments, retain camera,
recording/replay and settings controls, and need no telemetry-board connection.
See Section~\ref{sec:video} for exact receiver assignments.

\subsection{State-machine badge and telemetry meaning}
The large Systems badge follows the received state, including unexpected
transitions in recorded data. It never takes its value from a simulation.
\begin{longtable}{@{}L{.30\linewidth}L{.18\linewidth}L{.45\linewidth}@{}}
\toprule State & Color & Interpretation \\\midrule\endhead
\code{GROUND_TESTING} & Green & Received ground-testing state. \\
\code{PREFLIGHT} & Blue & Received preflight state. \\
\code{FLIGHT} & Red & Received flight state. \\
\code{POST_APOGEE}, \code{MAIN}, \code{END} & Orange & Recovery/post-apogee or end state. \\
Unknown or unavailable & Black & Missing telemetry or an unrecognized state. \\
\bottomrule
\end{longtable}
The caption identifies LIVE, DEMO or REPLAY. Stale LIVE reception retains the
last state with an explicit \ui{STALE} annotation. A state color does not
indicate command authorization, verified liftoff, motor burn or deployment.

\section{Connections, source isolation and command authority}\label{sec:connections}
\subsection{Local serial operation}
Launch has \textbf{one telemetry board}, used for downlink and intended
Launch-only uplink. Away has one downlink board and a separate pointer board.
Both existing serial protocols use 115200 baud, 8 data bits, no parity, one stop
bit, no flow control and the previous monitor's default DTR/RTS behavior.
Select the physical device and connect explicitly. Refresh after plugging in
equipment. A connected or connecting port disappears from the other board's
dropdown, including macOS tty/cu aliases, and returns after disconnection.

Connecting the telemetry board opens its port. \ui{Start Polling} separately
enables reception; \ui{Stop Polling} leaves the port open. The original ground
read timeout is one second; pointer reads use 0.1 second. No initialization
command is automatically sent on connection or polling start. Reconnect is
manual. Tracking never resumes automatically after stale data or disconnection.

\subsection{Computer connection versus rocket reception}
The link panel distinguishes disconnected, connected with polling stopped,
waiting for telemetry, live reception, and lost/stale reception. USB connected
means that a port opened. Rocket live means a checksum-valid LIVE packet arrived
within the mission freshness limit. There is no separate radio handshake.
Reconnection requires a new packet before reporting live radio reception.
Rejected packets do not update accepted telemetry or replay data.

\subsection{Launch uplink and Iris selector placeholders}
The Launch board's software uplink-enable command is not defined. LIVE rocket
commands therefore remain locked for Balius and Iris. The DEMO-only switch
simulates enable/disable and logs events. Away stations have no uplink authority.
Selecting a remote Away route or receiving a peer snapshot cannot unlock it.

For Iris the same Launch board has one Sustainer/Booster selection shared by
downlink and uplink. Away 1--3 initially desire Sustainer, Away 4 Booster.
The switching command and acknowledgement are also undefined. Selectors are
labelled \ui{PLACEHOLDER}; LIVE/REPLAY switching is unavailable and physical
target state is unknown. \ui{Simulate switch} in DEMO changes the indicated
simulated state without replacing or relabeling the Zephyrus recording.
Leaving DEMO clears this state. No fabricated selector or uplink-enable bytes
are transmitted.

\subsection{LIVE, DEMO and REPLAY}
LIVE starts disconnected. DEMO runs the recorded test flight and simulates
commands. REPLAY reads a saved session and cannot transmit. Mode changes close
transports, stop logging/tracking and clear incompatible history; DEMO and
REPLAY cannot construct physical serial workers. Virtual-pointer rehearsal is
an explicitly labelled port-free device in Away LIVE, independent of radio
readiness.
\begin{figure}[H]\centering
\begin{tikzpicture}
\node[accentbox] (live) at (0,0) {LIVE\\explicit connection\\to hardware};
\node[box] (demo) at (5.2,0) {DEMO\\recorded CSV playback};
\node[box] (replay) at (10.4,0) {REPLAY\\session opens paused};
\node[box,text width=126mm] (gate) at (5.2,-2.1) {Transition boundary: stop logging and tracking; close serial/video; invalidate old workers; clear visible history; install the selected source};
\draw[flow] (live)--node[above,font=\footnotesize]{select}(demo);
\draw[flow] (demo)--node[above,font=\footnotesize]{load}(replay);
\draw[flow] (live)--(gate.north west);
\draw[flow] (demo)--(gate);
\draw[flow] (replay)--(gate.north east);
\node[captionnode,text width=125mm] at (5.2,-3.45) {A planning-only flight returns to disconnected LIVE. A session-bearing flight restores paused REPLAY. Source-mode names are independent of the rocket's onboard flight state.};
\end{tikzpicture}
\caption{Mode isolation. Other mode selections use the same transition boundary; the upper arrows illustrate common operator paths rather than the only permitted transitions.}
\label{fig:modes}
\end{figure}


\section{Mission, locations and time}\label{sec:mission}
\subsection{Launch and antenna coordinate entry}
\ui{Mission $\to$ Configure} accepts launch latitude/longitude, MGRS or Plus
Code, all resolved offline to WGS84. MGRS accepts spaces/lowercase, UTM/polar
grids and displays precision; its decoded grid reference point is used.
Full Plus Codes resolve to the area's center. Short codes require nearby
reference latitude/longitude; locality names alone are not geocoded. Review
the resolved coordinates before saving. Switching display formats alone keeps
canonical coordinates; editing them changes a named preset to Custom.

Enter ellipsoid altitude for geographic transforms and MSL elevation for
weather/atmosphere separately, and establish the launch origin. The legacy
height interpretation may be unknown, ellipsoid GPS, GPS relative to launch,
or barometric relative to launch. Select the interpretation verified for the
firmware. This controls trajectory comparison, while physical legacy pointing
retains its separate original convention. Stop recording before mission edits.

At Away stations, antenna entry accepts lat/lon, MGRS or \ui{Relative to
launch}. Relative inputs are true heading \emph{from antenna toward the rocket
on its launch pad}, horizontal range 0--100,000 m, and antenna altitude minus
launch altitude. Heading 90 degrees, distance 500 m, height +5 m places the
antenna west of the pad and 5 m above it. Relative entry recalculates when the
launch origin changes; absolute entry retains its coordinates. Confirm
\ui{Antenna location established} for virtual trajectory tracking.

\Needspace{18\baselineskip}
\subsection{URRG presets}
\begin{longtable}{@{}L{.38\linewidth}L{.55\linewidth}@{}}
\toprule Location & MGRS \\\midrule\endhead
Rocket launch pad & \code{18TUN2061530290} \\
Launch station reference & \code{18TUN2063730181} \\
Away 1 antenna & \code{18TUN2177730106} \\
Away 2 antenna & \code{18TUN2291333237} \\
Away 3 antenna & \code{18TUN2015031101} \\
Away 4 antenna & \code{18TUN2259229678} \\
\bottomrule
\end{longtable}
Wizard selection fills pad and applicable Away antenna coordinates. The
Launch station's reference does not create a local pointer. Mission antenna
presets preserve the entered antenna altitude. None supplies surveyed elevation,
calibrates mechanics or sends a command. Saved missions/flights preserve their
own preset labels, exact codes and canonical coordinates.

\subsection{Time and flight alignment}
Device time is uptime, with rollover extension; it is distinct from host receive
time, recording elapsed time, UTC and simulation time. Flight zero is established
by an observed Preflight-to-Flight transition or explicit operator alignment.
Joining mid-flight leaves it unaligned. A detected reboot invalidates alignment;
realign deliberately. DEMO uses each receiver's own recorded timing. Replay
uses host elapsed time and restores alignment events. Video synchronization is
host timing with a configurable offset, not camera exposure synchronization.

\section{Rocket display, trajectory and GNC}\label{sec:flightview}
\subsection{Main and detailed telemetry}
The overview shows available altitude, velocity, GPS, radio/packet freshness,
phase and flight-time information. Systems adds GPS accuracy/satellite counts,
maximum altitudes, temperature, packet number, legacy servo drives and angles,
three battery cells, power/BMS and pyro values. Missing values stay unavailable.
GNC plots preserve the reported gyro rates and integrated rotations; new
canard/fin-tab feedback is not inferred from the four legacy servo channels.

Track coordinates are launch-relative east/north/up with an explicit origin
and altitude interpretation. A simulation reference is a separate dashed path
and same-time marker; it does not replace measurements. DEMO uses its own
labelled local display frame. Alerts and diagnostics expose accepted/rejected
packet counts, discarded bytes, inferred gaps, full samples and events.
Low-battery alerting uses a one-second dwell and clears 0.3 V above the
threshold. Acknowledgement records an event without erasing an active condition.

\subsection{Compact rocket orientation}
The large model replaces the circular attitude indicator; readouts and status
are underneath. The \ui{View azimuth} slider changes camera direction from
0 to 360 degrees while vertical stays vertical. This pane has no drag, pan,
scroll or hardware-control interaction.

\statusnote{Orientation source}{Zephyrus roll/pitch/yaw values are cumulative
body-axis gyro integrals, not Euler attitude. Their numerical readouts remain,
but the model is neutral and muted with \ui{ORIENTATION UNAVAILABLE}. Its upright
placeholder does not assert a measured vertical attitude. Flight phase or
climbing altitude cannot establish orientation.}

A future adapter can supply an explicitly validated body-to-ENU quaternion.
Current telemetry also lacks confirmed motor-burning and parachute-deployed
fields, so those effects are unknown. Pyro fired and flight state do not prove
deployment; historical ignition does not prove continued burning.

Choose \ui{OpenRocket reference} explicitly for simulated orientation and
events. This focuses the embedded Launch 3D controls or opens the Away 3D
window. Playback, scrub and event jumps drive the reference pose. The pane
labels reference time and playback state. A simulated burn interval displays a
flame; a recorded simulated deployment displays a canopy. Switching back to
Telemetry restores its independent values and availability.
\begin{figure}[H]\centering
\begin{tikzpicture}
\node[box,text width=39mm] (raw) at (0,0) {Current Zephyrus\\body-gyro integrals};
\node[box,text width=39mm] (valid) at (5.3,0) {Future adapter\\validated body-to-ENU pose};
\node[box,text width=39mm] (ref) at (10.6,0) {Selected OpenRocket\\quaternions + events};
\node[box,text width=39mm] (neutral) at (0,-2.1) {Raw numeric readouts\\neutral model\\orientation unavailable};
\node[accentbox,text width=39mm] (pose) at (5.3,-2.1) {Telemetry pose/effects\\only with explicit valid data};
\node[accentbox,text width=39mm] (sim) at (10.6,-2.1) {Reference pose/effects\\labelled simulation time};
\draw[flow] (raw)--(neutral);
\draw[flow] (valid)--(pose);
\draw[flow] (ref)--(sim);
\node[captionnode,text width=140mm] at (5.3,-3.65) {No cross-filling: simulation does not repair missing telemetry.\quad Camera azimuth changes viewing direction only.};
\end{tikzpicture}
\caption{Orientation provenance. Phase, increasing altitude and pyro bits do not supply missing orientation, motor-burning or parachute-deployment evidence.}
\end{figure}


\subsection{Full 3D flight viewer}
Use \ui{View $\to$ 3D flight} after loading or simulating a reference. Drag to
orbit, Shift/right-drag to pan, scroll to zoom and double-click or Fit to reset.
Perspective/Top/Side views, follow-rocket, display-size, velocity and ground
projection controls affect viewing only. Play, rewind, speed, timeline and
event jumps inspect ignition, burnout, apogee, deployment and touchdown when
those events exist. \ui{Follow monitor time} uses aligned flight time or
virtual rehearsal time; turn it off for independent playback.

Gold denotes reference, teal actual track when origins agree, purple antenna.
Old four-column references remain usable with an explicitly path-aligned
illustration and unavailable attitude/events. Tumble/recovery can leave
engine-held attitude; undersampled rotation is labelled. Missing events are
never inferred from descending motion. The rocket, plume and canopy are
schematics, not CAD or deployment dynamics. The Zephyrus model has no recovery
event in its nominal output; the bundled small example supplies one.

\section{Away antenna pointer}\label{sec:antenna}
\subsection{Model and manual operation}
The four-leg stand supports an azimuth turntable and an elevation carrier.
Grid dish, long Yagi and enclosed Avenger XR18 attach to the common beam
through the support pivots. The dish-colored feed arm is collinear with the
boresight. Dimensions and hidden mechanics are illustrative. Drag to orbit
from any angle, scroll to zoom, and double-click or choose a preset to reset.
Camera movement is independent of command authorization.
\begin{figure}[H]\centering
\begin{tikzpicture}
\node[box,text width=35mm] (stand) at (0,0) {Four-leg stand\\fixed world frame};
\node[box,text width=35mm] (turn) at (5,0) {Turntable and supports\\azimuth rotation};
\node[accentbox,text width=35mm] (beam) at (10,0) {Common pivot beam\\azimuth and elevation};
\node[smallbox] (dish) at (5.7,-2) {Grid dish\\and aligned feed};
\node[smallbox] (yagi) at (10,-2) {Long Yagi\\beam attachment};
\node[smallbox] (helix) at (10,-3.8) {Avenger XR18\\beam attachment};
\draw[flow] (stand)--(turn);
\draw[flow] (turn)--(beam);
\draw[flow] (beam.south west)--(dish.north);
\draw[flow] (beam)--(yagi);
\draw[flow] (beam.east) -- ++(.6,0) |- (helix.east);
\node[captionnode,text width=42mm] at (.3,-2.8) {Orbit and zoom change the camera only. They do not change commanded azimuth or elevation.};
\end{tikzpicture}
\caption{Transform hierarchy of the procedural antenna model. Both outer antennas and the reflector belong to the same moving carrier; this diagram describes attachments, not dimensions or a photographic viewing direction.}
\label{fig:mount}
\end{figure}


Enter azimuth/elevation and press Send. Editing a field alone does not move or
transmit. UP/DOWN/LEFT/RIGHT send the original direction opcodes; ZERO resets
software counts at the physical current position, rather than homing to a
switch. Manual controls require only the pointer connection, independently of
telemetry. A direction command before a known target leaves angles unknown.

\ui{Sent} means serial dispatch, not acknowledgement or measured arrival.
The graphic smoothly approaches the last target on its independent animation
timer. Smoothing changes no commands or timing. Measured pose is unavailable
in the existing pointer protocol.

\subsection{Physical tracking}
Connect both boards, start polling and use \ui{Send to AntPtr} to freeze ground
receiver GPS. \ui{Start tracking} sends targets at up to 5 Hz using rocket
GPS lat/lon and filtered barometric height. It preserves receiver lat/lon
rounded to five decimals, receiver height fixed to zero for pointing, and the
historical displayed receiver-height division by 1000. Mission trajectory
datum settings do not alter this legacy path. Unfreeze updates receiver position.
Stale/invalid telemetry or disconnection stops host tracking; restart explicitly.

\statusnote{Stopping host tracking}{Stop tracking removes unsent work and stops
new targets. The last commanded physical move may continue. The firmware has
no stop/watchdog, measured pose, mechanical homing or acknowledgement. Use the
equipment's independent stop/power control to interrupt physical movement.}

\subsection{Virtual pointer and trajectory rehearsal}
In Away LIVE, configure an established antenna location and load a georeferenced
reference. Select \ui{Virtual antenna pointer}, connect, and use the same
manual commands or \ui{Follow trajectory}. No telemetry connection is needed.
The simulation slews at illustrative maximum rates of 90 degrees/s azimuth
and 60 degrees/s elevation. Playback has pause/resume, rewind, time scrub and
0.25--4 times speed. Manual commands pause rehearsal; Hold freezes its clock
and motion. Changing reference/location stops it and requires restart.

The banner, pose and connection are labelled VIRTUAL/SIMULATED. A gold target
and purple antenna marker remain distinct from telemetry. The stored reference
origin and antenna ellipsoid position define the direction. Coincident points
have no direction; exact zenith retains the previous azimuth. Rehearsal neither
opens a serial port nor creates telemetry or unlocks rocket controls. Changing
source mode closes it; flight loads do not restore a simulated connection.
\begin{figure}[H]\centering
\begin{tikzpicture}
\node[box,text width=36mm] (reference) at (0,0) {OpenRocket reference\\ENU path + launch origin};
\node[accentbox,text width=36mm] (frame) at (5.15,0) {Antenna-local direction\\ENU via ECEF};
\node[box,text width=36mm] (pose) at (10.3,0) {Virtual pointer\\command + slew model};
\node[box,text width=36mm] (mission) at (0,-2) {Mission profile\\antenna lat/lon/height};
\node[box,text width=36mm] (clock) at (5.15,-2) {Independent clock\\play, pause, seek, speed};
\node[box,text width=36mm] (display) at (10.3,-2) {Simulated pose + markers\\telemetry unchanged};
\draw[flow] (reference) -- (frame);
\draw[flow] (frame) -- (pose);
\draw[flow] (mission.east) -- (frame.south west);
\draw[flow] (clock) -- (frame);
\draw[flow] (pose) -- (display);
\end{tikzpicture}
\caption{Away virtual rehearsal uses the saved antenna position and its own clock. No physical serial transport participates in this path.}
\end{figure}


\section{Video roles and routing}\label{sec:video}
Digital/Analog labels describe the rocket-to-ground RF paths. All locally
selected inputs are USB camera captures. Telemetry is not required to start
video or a video-only recording.

\begin{longtable}{@{}L{.21\linewidth}L{.31\linewidth}L{.41\linewidth}@{}}
\toprule Profile & Local receivers & Large Launch channels / source assignment \\\midrule\endhead
Balius Away 1--4 & Digital + Analog & Every Away contributes one of each. \\
Balius Launch & Digital only & Digital initially local; Analog initially Away 1. \\
Iris Away 1--3 & Sustainer Digital + Sustainer Analog & Sustainer telemetry/APRS assignment. \\
Iris Away 4 & Sustainer Digital + Booster Analog & Booster telemetry/APRS assignment. \\
Iris Launch & Sustainer Digital only & Digital initially local; Sustainer Analog initially Away 1; Booster Analog initially Away 4. \\
\bottomrule
\end{longtable}
Iris has three system channels but \emph{two local USB feeds per Away station}.
The later receiver-assignment diagram refines the earlier three-local-camera
proposal. Stage names are assignments, not decoded identification from legacy
packets. Launch's one telemetry board targets one stage at a time.
\begin{figure}[H]\centering
\begin{tikzpicture}
\node[box,text width=38mm,minimum height=20mm] (launch) at (0,0) {\textbf{Launch USB}\\Sustainer Digital};
\node[box,text width=38mm,minimum height=20mm] (away) at (5.3,0) {\textbf{Away 1, 2, 3}\\Sustainer Digital\\Sustainer Analog};
\node[box,text width=38mm,minimum height=20mm] (four) at (10.6,0) {\textbf{Away 4}\\Sustainer Digital\\Booster Analog};
\node[accentbox,text width=144mm,minimum height=20mm] (large) at (5.3,-3) {\textbf{Launch primary views: Iris}\\Sustainer Digital\qquad Sustainer Analog\qquad Booster Analog};
\draw[flow] (launch)--(large.north west);
\draw[optional] (away)--node[right,font=\footnotesize]{future transport}(large.north);
\draw[optional] (four)--(large.north east);
\node[box,text width=144mm] (thumbs) at (5.3,-5) {Eight Away thumbnails: four station pairs remain visible\\Double-click replaces only the matching primary channel};
\draw[flow] (thumbs)--(large);
\end{tikzpicture}
\caption{Iris has three system channels and two local inputs per Away station. For Balius every Away pair is Digital/Analog and Launch has two primary views. Dashed paths are not implemented network feeds.}
\label{fig:video}
\end{figure}


\subsection{Capture and device exclusivity}
For each input choose Find, select a USB receiver and Start. Stop closes only
that channel; File selects local video. Each panel reports frame age and faults.
Grant OS camera permission when required. A selected camera disappears from
other local selectors; active and closing cameras remain reserved. Clear an
unused selection to release it. One failed/reconfigured stream does not stop
another. These reservations apply within the application's local controllers,
not as a cross-process camera arbitration service.

\subsection{Launch video wall}
All four Away pairs remain visible as thumbnails, including lower-quality
alternatives. Double-click one to replace \emph{only its matching large
channel}. Away 4 Booster Analog replaces Booster Analog; Digital always
replaces Digital. \ui{Use local Digital} restores the Launch USB feed without
having closed its capture during promotion.

Remote slots currently show \ui{Awaiting station link}. Choosing a source does
not establish a link. The presentation interface accepts future frames,
quality text, receive age and loss/stale states; five seconds without fresh
remote frames becomes stale. Interstation video transport and remote recording
are not implemented. The existing telemetry LAN link carries no video.

\subsection{Recording and replay}
After local frames arrive, Start recording creates a session independently of
serial telemetry. Channels have separate FFmpeg segments and timing. DEMO
video is explicitly a TEST PATTERN, not Zephyrus flight footage. REPLAY exposes
flight/session load, play/pause, seek and speed in the simplified footer.
Launch replays its local Digital channel; inactive analog assets remain in
the archive. Remote frames are display-only in the present integration API.

\section{Legacy rocket commands, power and recovery}\label{sec:rocket}
Systems retains Advance State, Zero P/Y/Roll, Zero Altitude, Zero Velocity,
Zero Servos, PD Activate, roll/airbrake angle entry and VTX 1/3/5/8 W. Six
power-rail requests/readbacks cover 3, 3.3, 5, 7.4, 8.4 and 28 V, with current,
voltage, total and temperature. The 3.3 V request remains on; Total is an
aggregate indicator. BMS displays all eight enabled/triggered protection bits.

Recovery has six channels with continuity, arm/fire status and resistance.
Select channels 0--5 for ARM/FIRE. State advance, ARM/FIRE and emergency piston,
BP wells, tender descender and all-recovery actions retain their confirmation
dialogs and original packet encoding. These controls are present for parity;
normal Launch LIVE dispatch remains locked pending the uplink switch contract.
In DEMO they can be exercised after simulated uplink enable. Neither a command
event nor a fired bit proves a physical recovery outcome.

\section{Wind and nominal OpenRocket reference}\label{sec:simulation}
\subsection{Wind preparation}
Enter/import layers or retrieve online weather for a specified site and a
timezone-bearing time, e.g. \code{2026-10-25T16:00:00Z}. Apply after table edits.
Heights are above launch in metres, speed is m/s and direction is meteorological
\emph{from} true north. Internal layers store east/north velocity components.
Layers must have increasing nonnegative height; interpolation holds endpoint
wind beyond coverage.

The weather adapter uses the nearest available hour within one hour and
converts geopotential height to geometric MSL height, then subtracts launch
MSL elevation. Review available altitude/time coverage. Requests/responses are
retained for offline reuse. Weather is attributed to
\href{https://open-meteo.com/en/docs}{Open-Meteo}, CC BY 4.0. Manual profiles
remain available without network access.

\subsection{Simulation and import}
Select an ORK with a saved simulation/motor configuration. Simulation index
zero selects the first saved simulation. The bundled MITRT N8406 curve resolves
the supplied Zephyrus test-launch model; other custom ENG/RSE curves can be
added. Missing motors are reported by name. Run simulation in a separate job;
Cancel stops it without interrupting telemetry. Inspect the job's request,
copied inputs, hashes, trajectory metadata and worker log on failure.

The bridge uses the selected saved motor/flight configuration, configured
site/rail/wind, ISA atmosphere, fixed seed, zero turbulence standard deviation
and first output branch. It disables saved extensions and the fork's global
inertia override. Branch count and warnings are recorded. Controlled
three-axis canard/tab dynamics remain unqualified; this is a nominal reference.

Imported references use CSV \code{time_s,east_m,north_m,up_m} with matching
JSON declaring schema 1, ENU, units \code{m,s} and geographic origin
\code{[latitude,longitude,ellipsoid_altitude]}. Up is launch-relative.
LIVE requires compatible origin and nonsynthetic data. References are not
extrapolated outside their time coverage. Optional quaternion, rate, velocity
and event data support 3D playback; old position-only references remain usable.

\section{Recorded test-launch demonstration}\label{sec:demo}
Select DEMO explicitly. GS2 is the default; GS1 and GS3 are independently
selectable. The bundled gzip files reproduce the supplied CSV bytes offline.
These are three receivers of the same launch, not successive flight segments
and not a fused dataset.
\begin{longtable}{@{}L{.23\linewidth}L{.22\linewidth}L{.25\linewidth}L{.23\linewidth}@{}}
\toprule Receiver & Rows & Receive duration & Notes \\\midrule\endhead
GS1 & 4,366 & 274.056 s & GPS corruption present. \\
GS2 & 5,464 & 330.234 s & Default, greatest coverage. \\
GS3 & 4,939 & 332.725 s & No valid receiver GPS. \\
\bottomrule
\end{longtable}
Playback starts five receive-time seconds before the first FLIGHT packet,
device time 747.766 s. Play/pause, seek, 0.25--4 times speed, Launch $-5$ s and
Full start control playback. Gaps, original UTC/packet numbers and values are
preserved. Seeking reconstructs prior history. Stop logging before changing
recording or seeking; pause/speed changes are allowed during logging.

All logs end during MAIN descent, with no touchdown. They include an
approximately 8.4-second reception gap, absent fixes, zero cell readings and
gyro saturation. GS3 briefly reports Preflight during descent. The app retains
these values, rather than correcting flight history. Track-only checks omit
missing 3D fixes and points over 20 km from launch; the banner identifies them.
The display uses GPS horizontal offsets and reported barometric height, with
no surveyed-datum claim. GS3 cannot rehearse automatic receiver-GPS tracking;
GS1/GS2 can. Manual simulated controls remain available.

No raw frames or launch video were supplied. CSV checksums cannot be
revalidated; demo video is a labelled test pattern. Commands are simulated.
Returning to LIVE clears demo data and restores the prior mission/reference.

\section{Recording, sessions and portable flights}\label{sec:persistence}
\subsection{Capture and session folders}
Start Logging immediately creates a unique folder when eligible telemetry or
video is available. The default is the app-data sessions folder, overridable
in Settings. Diagnostics shows its path. The recorder is independent of
redraws; its drop/error/completion status is separate from UI frame drops.
Stop Logging finalizes outputs; quitting also closes active workers/recording.

Each session may contain a manifest, indexed SQLite samples/events, CRC-protected
raw serial RX/TX, 43-column \code{telemetry.csv} and rejected-packet CSV, a
reference, imported source CSV and per-channel video segments. Opening a session
supports play/pause, seek, 0.1-second step and 0.25--4 times speed. Seek removes
future state and restores prior alignment and sent-pointer events. CSV export
of normalized samples includes full sample JSON and differs from legacy CSV.

Video follows host timing. Positive mission offset delays it; negative advances
it. Paused seeking decodes one frame per recorded stream. Stopping a replay
feed keeps it paused until the next seek/restart. Old single-stream recordings
map to Digital. Reconnect creates a new segment directory. Incomplete sessions
can be read when committed SQLite data survives; retain a crashed session's
\code{-wal} file for recovery. A crash can lose the active segment/uncommitted
samples without invalidating finalized data.

\subsection{One-file save/load}
\ui{File $\to$ Save flight} writes a portable \code{.rktflight}; Open flight
restores it. The file includes mission/site/wind/weather, model/custom motors,
reference and available recordings. Embedded assets no longer depend on
their original paths. Saving is atomic; loading validates inventory, paths,
hashes and contents before replacing current state.

\begin{longtable}{@{}L{.29\linewidth}L{.64\linewidth}@{}}
\toprule Save context & Content and limit \\\midrule\endhead
Active logging & Committed snapshot; finalized video segments only. Stop Logging and save again for the closing segment and final samples. \\
Stopped logging & Last session in the current source mode; waits for video closure. \\
REPLAY & Opened session and playback position. \\
DEMO without recording & Entire selected receiver dataset and current position, with no invented raw packets or flight video. \\
LIVE without recording & Displayed buffer, at most 6,000 samples, labelled with that limited scope. \\
No telemetry & Planning flight with mission/model/reference. \\
\bottomrule
\end{longtable}
Stop logging before opening another flight. A session-bearing file opens paused
in REPLAY with hardware disconnected. A planning file opens disconnected LIVE.
Pointer calibration and connections are not restored. Mission JSON separately
stores configuration and can retain external asset paths; use a flight file
for portable assets. Extracted flights live in the app-data flight cache.

Legacy CSV import accepts the audited layout, retains original bytes/extra
columns and converts to replay. Other CSV layouts are rejected. Imported
decoded rows cannot recover radio checksums or establish new vehicle feedback.

\section{Station networking and fusion status}\label{sec:network}
Launch's intended topology is Ethernet/PoE to four LTU-XR links, each ending at
an Away station with its local WLAN. Away Wi-Fi selection stores a desired SSID;
Refresh reads OS current/saved networks and Open Wi-Fi settings opens the OS
joining controls. The selected SSID and actual OS connection are distinct.
The monitor stores no password and does not join or configure the bridges.

Launch's Away pointer selector records a desired communication route. It does
not move a pointer, enable uplink or prove a link. Remote telemetry/video/control
transport, LTU signal metrics, analog receiver tuning and interstation calls
are pending. No sensor fusion is implemented: local telemetry, peer snapshots,
APRS and simulation remain separate sources.

\subsection{Existing one-peer LAN}
Join the local network, Listen on one computer (default TCP 8765), then Connect
from the other to its LAN IP and port. The listener binds local IPv4 interfaces
and accepts one peer. About twice per second the two sides exchange bounded,
versioned, read-only snapshots of selected telemetry, pointer/link state and
available APRS. The panel shows peer name/mode/link, altitude, phase and age.
Stop closes the link and removes disconnected peer data. This connection
carries no video/audio/commands and does not supply the future four-station
transport or a fused navigation solution.

\subsection{APRS recovery information}
Configure an external Dire Wolf audio decoder and KISS TCP, then connect the
APRS row (default \code{127.0.0.1:8001}). An optional full callsign including
SSID filters fixes; blank accepts the latest supported position. The panel
reports connection, callsign, coordinates, optional altitude and age. A
disconnected last fix stays labelled as such.

Supported AX.25 UI payloads are uncompressed APRS \code{!}, \code{=},
\code{/} and \code{@} positions. Compressed, Mic-E, objects, NMEA and ambiguous
positions are unsupported. \code{/A=} altitude converts feet to metres without
assuming a datum. APRS is receive-only recovery information and cannot replace
telemetry or actuate pointing. Automatic Iris identification/relay is pending.
Host/port/filter settings persist in \code{station_network.json}; connections
start off. Dire Wolf is neither bundled nor configured by this app.

\section{Settings, shortcuts and practical checks}\label{sec:settings}
\subsection{Global Settings}
Open Settings with Command-comma on Mac or Ctrl-comma on Windows. The dialog is
available disconnected and in every role. OpenRocket JAR selects a compatible
complete team engine; Use bundled restores the packaged engine. Its effective
path/version are displayed. Simulation timeout is 10--1800 s, default 120.
Recording folder affects future sessions only. Dark/Daylight appearance
persists; the View menu toggle shares that preference.

Save validates and stores; Cancel discards edits. Restore Defaults updates the
form and takes effect on Save. \code{settings.json} resides in the displayed
per-user data directory; \code{--data-dir} isolates it. Global preferences are
not included in mission or flight files. Running simulations keep their
submitted engine/timeout. Missing custom JARs produce an error rather than
silently falling back. Startup never auto-connects equipment.

\Needspace{22\baselineskip}
\subsection{Preserved and added keyboard shortcuts}
Qt maps Ctrl to Command and Alt to Option on macOS, as in the original UI.
Actions retain role/connection gating; shortcuts cannot bypass a disabled
hardware capability. Launch has no pointer shortcuts to a local device; Video
hides serial actions. The ten original serial/pointer key sequences are:
\begin{longtable}{@{}L{.36\linewidth}L{.28\linewidth}L{.29\linewidth}@{}}
\toprule Action & Windows / Qt & macOS \\\midrule\endhead
Refresh serial ports & Ctrl+R & Command+R \\
Telemetry connection toggle & Ctrl+Alt+C & Command+Option+C \\
Pointer connection toggle & Shift+Ctrl+Alt+C & Shift+Command+Option+C \\
Start/stop polling & Ctrl+Return & Command+Return \\
Start/stop logging & Ctrl+L & Command+L \\
Pointer UP/DOWN/RIGHT/LEFT & Ctrl+arrow & Command+arrow \\
Pointer ZERO & Ctrl+0 & Command+0 \\
\bottomrule
\end{longtable}
Open/Save flight add Ctrl+O/Ctrl+S (Command on Mac); Settings adds Ctrl+comma
(Command on Mac). Forget setup deliberately has no shortcut. Serial actions
are bound once and work from either Launch display.

\subsection{Operating sequence}
Choose station/role/vehicle/site; review elevations and mission height datum.
Connect the appropriate local inputs explicitly. For telemetry, start polling
and distinguish USB readiness from recent rocket reception. For video, check
each channel label and independent input. Start logging before the interval
that must be retained. On Away, verify pointer orientation and ground-GPS freeze
before physical tracking, or use explicitly virtual rehearsal.

After the flight, stop host tracking, stop logging and inspect completion/error
status. Save a final flight, reopen it paused and confirm expected telemetry,
reference and local video. Retain the original session until the archive has
been checked. Camera permissions, actual board I/O, pointer mechanics and the
station network require equipment checks beyond the software test suite.

\begin{longtable}{@{}L{.31\linewidth}L{.62\linewidth}@{}}
\toprule Symptom & Check \\\midrule\endhead
USB connected, no rocket data & Start Polling; verify board/port and radio. A port open is not a radio handshake. Inspect rejected packets and freshness. \\
Commands locked at Launch & Expected in LIVE until the software uplink-enable contract is implemented. DEMO exercises its labelled simulation. \\
Camera absent from selector & It may be selected, active or closing in another channel. Release it there, refresh, and check OS permission. \\
Rocket upright but unavailable & Neutral placeholder: gyro integrals do not establish attitude. Choose a valid reference explicitly for simulated pose. \\
No flame/canopy & Current telemetry has no confirmed fields; a reference may lack the required events. Do not infer them from phase. \\
Simulation has no usable path & Verify active saved motor/configuration, custom curves, chosen JAR and worker log; fewer than two valid rows is rejected. \\
Tracking stopped & Inspect freshness, GPS and connection; restart deliberately. A host stop cannot stop the last physical move. \\
Remote thumbnails empty & Station video transport is pending; source selection alone does not create a connection. \\
Controls crowded & Use the intended Launch display area; small-window layout improvements remain outstanding. \\
\bottomrule
\end{longtable}
